\documentclass{these-ISSS} % Pour la version finale
%\documentclass[draft]{these-ISSS} % Pour une version brouillon
%\documentclass[gray]{these-ISSS} % Pour une version en grayscale

\usepackage{mesmacros}
\usepackage{import}
\usepackage{apacite}
\usepackage{lipsum}
\usepackage{dsfont}
\usepackage{booktabs}
\usepackage{multirow}
\usepackage{footnotehyper}
\setlength{\headheight}{13.9186pt}
\renewcommand{\thefootnote}{\arabic{footnote}}
\renewcommand{\thempfootnote}{\arabic{mpfootnote}}
\newcommand{\indicator}{\mathds{1}}

\def\dataset{\ensuremath{\mathcal{D}}}
\def\x{{\mathbf{x}}}
\def\y{{\mathbf{y}}}
\def\b{{\mathbf{b}}}
\def\X{{\mathbf{X}}}
\def\o{{\mathbf{o}}}
\def\z{{\mathbf{z}}}
\def\C{\ensuremath{\mathcal{C}}}
\def\real{\mathbb{R}}
\def\interger{\mathbb{N}}
\def\E{\mathbb{E}}
\newcommand{\prob}{\mathbb{P}}
\def\localupdaterule{\mathrm{Local\_Update}}
\def\dataset{\ensuremath{\mathcal{D}}}
\def\localdataset{\dataset_u}
\def\subclients{\ensuremath{\mathcal{U}_s}}
\def\clients{\ensuremath{\mathcal{U}}}
\def\model{\theta}
\def\samplespace{{\mathcal{Z}}}
\def\messages{\ensuremath{\mathcal{M}}}
\def\roundsclient{\ensuremath{\mathcal{T}_u}}
\def\mimicfunction{{\obs}}
\def\model{\theta}
\def\localmodel{\theta_u^*}
\def\localmodelsen{\theta_u^{*,s}}
\def\localmodelpub{\theta_u^{*,p}}
\def\globallossfunction{{\loss}}
\def\locallossfunction{{\loss_u}}
\def\weights{w}
\def\gradient{\mathbf{g}}
\def\cc{\mathbf{u}}
\def\locallossfunctioncc{{\loss_\cc}}
\def\locallossfunctionbc{{\loss_{\bar u}}}
\def\localmodelcc{\theta_\cc^*}
\def\y{\mathbf{y}_u}
\def\xhalf{\mathbf{Z}_u}
\def\H{\mathbf{H}}
\def\i{\mathbf{I}}
\def\W{\mathbf{W}}
\def\V{\mathbf{V}}
\def\vv{\mathbf{v}}
\def\a{\mathbf{A}}
\def\u{\mathbf{U}}
\def\g{\mathbf{g}}
\def\trace{\mathbf{trace}}
\def\e{\mathbf{\epsilon}}
\def\ga{\mathbf{\gamma}}
\def\var{\mathbb{V}\mathrm{ar}}
\def\distribution{\ensuremath{\mathcal{D}}}
\def\extacc{Acc^{Ext}_u}
\def\acc{Acc_u}
\def\algo{\ensuremath{\mathcal{A}}}
\def\s{\mathbf{s}_u}
\def\hats{\hat \s}
\def\ham{d_{H}(\s,\hats)}
\def\P{\mathbf{P}}
\def\ei{\mathbf{e}(\model)}
\def\e{\mathbf{e}}
\def\h{\mathbf{h}}
\def\prior{\rho}
\def\truepositive{\ensuremath{\mathcal{P}}}
\def\truenegative{\ensuremath{\mathcal{N}}}
\def\domain{\mathbb{D}}
\def\subtime{\mathbf{T}}
\def\Win{\W^{a}}
\def\Wout{\W^{L}}
\def\bin{\b^{a}}
\def\bout{\b^{L}}
\def\oi{\mathbf{g}_{i}}
\def\oj{\mathbf{g}_{j}}
\def\ob{g_{b_{i}}}
\def\intervals{\ensuremath{\mathcal{I}}}
\def\obs{\ensuremath{\mathcal{G}}}
\def\bini{b_i^{a}}
\def\obsnew{\obs_{\mathrm{new}}}
\def\loss{\ensuremath{\mathcal{L}}}
\def\alphas{\mathbf{\alpha}}
\def\batch{\ensuremath{\mathcal{B}}}
\def\zini{z^{a}_i}
\def\gradlossjwi{\frac{\partial \loss_j}{\partial \W_i^a}}
\def\gradlossjbi{\frac{\partial \loss_j}{\partial b_i^a}}
\def\gradlossjW{\frac{\partial \loss_j}{\partial \W}}
\def\gradlossjb{\frac{\partial \loss_j}{\partial \b}}
\def\gradlosswi{\frac{\partial \loss}{\partial \W_i}}
\def\gradlossbi{\frac{\partial \loss}{\partial b_i^a}}
\def\gradlossjzi{\frac{\partial \loss_j}{\partial z_i}}
\def\gradzibi{\frac{\partial z_i}{\partial b_i^a}}
\def\activeset{\mathcal{A}}
\def\gradlosswiatt{\frac{\partial \loss}{\partial \W_i^a}}
\def\gradlossbiatt{\frac{\partial \loss}{\partial b_i^a}}
\def\slice{\mathbf{s}}
\def\gradziwi{\frac{\partial z_i}{\partial \W_i^a}}
\DeclareMathOperator*{\argmax}{argmax}
\DeclareMathOperator{\Dist}{\textbf {Dist}}

\newcommand{\francesco}[1]{\textcolor{red}{\textbf{Francesco:} #1}}


% Redefine end statements to nothing

\newcommand{\eps}{\ensuremath{\varepsilon}}
\newcommand{\relu}{\textrm{ReLU}}

\DeclareMathOperator*{\argmin}{argmin}

% Biblio pour les pages webs
\newcites{web}{Pages web}
% Biblio pour mes publications
\newcites{mine}{Mes publications}

\def\mytt#1{\small\texttt{#1}}

\input{titreEtInfos}


\begin{document}
  \pagenumbering{roman}

  \input{jury}

  \maketitle

  \begin{dedicace}
    À toi lecteur <3 :*
  \end{dedicace}

  \cleardoublepage

  \input{resume}

  \input{remerciements}

  \tableofcontents

  \mainmatter

  % \chapter{Introduction}\label{chap:intro}
% \section{Context of the Thesis}\label{sec:context}
% In the last years, the increasing availability of connected devices has progressively shifted the data generation process from centralized infrastructures towards large populations of edge devices, including smartphones, wearables, vehicles and sensors.
% According to~\citeA{cisco2023annual}, the number of networked devices and connections was expected to reach 29.3 billion by 2023, corresponding to more than three times the global population, and approximately 3.6 devices per person. Each of these devices represents a valuable data generation source, collecting information about daily interactions between individuals, devices and digital services.

% In this context, Machine Learning (ML) emerges as an incredibly powerful technology to automatically extract patterns from enormous amount of data. In general, access to more data allows to capture richer statistical regularities, improve generalization, and adapt to a wider range of users, environments and tasks~\cite{banko-brill-2001-scaling, shalev2014understanding, hestness2017deeplearningscalingpredictable, kaplan2020scalinglawsneurallanguage}. 

% However, having access to more data does not necessarily imply that these data can be freely collected or used. In many domains, data are subject to legal and ethical constrints~\cite{eu_gdpr_2016, hiroshima_process_guiding_principles_2023, eu_ai_act_2024, Buick_copyright, eu_ehds_2025, california_ccpa_cpra}. Medical records are protected because they describe a person’s health status. Financial data may reveal income, consumption patterns, debt, or economic vulnerability. Text, images, and user-generated content may be protected by copyright or by terms of service. Location traces may disclose where a person lives, works, studies, or receives medical care. Even apparently harmless digital traces can become sensitive when aggregated over time or combined with other sources of information \cite{netflix_attack}. Data are therefore not only an input to optimization procedures but are also representations of individuals, their habits, preferences, behaviors, and social relations. 

% This observation raises a broader question about control. If data generated by personal devices can reveal intimate aspects of people’s lives, who has access to such data, and under which conditions? In practice, large private companies often occupy a privileged position in the digital ecosystem. They design the devices, operating systems, applications, cloud services, and platforms through which data are generated, stored, and processed. As a consequence, they may acquire detailed knowledge about users’ behavior, preferences, movements, purchases, communications, and interactions. This knowledge can be used in many useful ways to improve people’s daily living, but at the same time, it can also be used for profiling, targeting, ranking, persuasion, facial recognition, or other forms of automated decision-making that are not always transparent to the individuals concerned and are beneficial only for few peoples \cite{pasquale_black_box_society_2015, veliz_privacy_is_power_2021}.

% The question is therefore not only whether data are useful, but whether the entities that collect and process them can be trusted. Recent enforcement actions involving major technology companies suggest that privacy risks often arise not only from external attacks or accidental data breaches, but also from ordinary platform practices.

% The Cambridge Analytica scandal showed how personal information collected through a digital platform could be repurposed for political profiling and targeted persuasion, raising public concern about consent, transparency, and the secondary use of personal data \cite{ico_democracy_disrupted_2018, ico_data_analytics_political_campaigns_2018}. This episode is part of a broader pattern of documented privacy controversies involving major technology companies \cite{cnil_google_gdpr_2019, ftc_youtube_coppa_2019, ftc_twitter_security_data_ads_2022, ftc_amazon_alexa_coppa_2023, texas_ag_google_privacy_2025}. These cases suggest that privacy risks do not arise only from external breaches, but also from the ordinary ways in which platform operators collect, retain, combine, and repurpose personal data.

% Similar concerns arise in machine learning, where the demand for large and diverse training datasets creates strong incentives to collect and reuse data at scale. Controversies involving biometric systems trained on scraped or improperly obtained facial images \citeweb{noyb_clearview_criminal_complaint_2025}, copyrighted material used for model training \citeweb{reisner_ai_pirated_books_problem_2025}, and wearable devices capturing data from both users and bystanders \citeweb{svd_meta_ai_glasses_privacy_2026} illustrate how the search for training data can blur the boundary between legitimate data use and unauthorized appropriation.


% These concerns highlight two obstacles to the conventional centralized learning paradigm. The first is practical. Data are increasingly generated across large populations of devices, so collecting them into a single repository may require substantial communication and storage costs, making centralized collection inefficient and difficult to maintain at scale. 

% The second obstacle concerns privacy and trust. Even when centralization is technically feasible, transferring raw records to a server increases the exposure of sensitive information and requires users or institutions to trust the entity that collects, stores, and processes the data. The problem is therefore not only how to learn from more data, but how to learn from distributed data without requiring all raw records to be pooled under the control of a single actor.

% \section{Federated Learning} \label{sec:fed_learning}

% Federated Learning (FL) has emerged as a response to this challenge \cite{konecny2015federatedoptimizationdistributedoptimizationdatacenter, pmlr-v54-mcmahan17a}. Instead of collecting raw data in a central repository, FL moves the training procedure closer to the data sources. Clients keep their datasets locally, compute model updates on their own devices, and communicate only these updates to a central server, whose role is to coordinate the training process and aggregate the received information. In the federated setting, the information exchanged with the server is limited, at least by design, to the updates required for collaborative model training.

% This design addresses two aspects of the centralized training. First, it reduces the practical burden of centralized data collection, since transmitting model updates can be substantially more efficient than transferring large volumes of raw data to a central storage infrastructure. Second, by avoiding the direct transmission of raw records, FL is naturally aligned with the principle of \textit{data minimization}, according to which only data that are directly relevant and necessary to accomplish a specified purpose should be collected \cite{eu_gdpr_2016, eu_regulation_2018_1725, TRUONG2021102402}.

% This second property has made FL particularly attractive in privacy-sensitive domains, including healthcare and finance \cite{Rieke2020, ZHANG2021106775, 10.1145/3298981}. In such settings, FL appears to support the idea of \textit{privacy by design}: privacy considerations are incorporated into the learning architecture itself, rather than being treated only as an external constraint. This intuition is also reflected in the original formulation of FL, where the authors observe that the updates exchanged during training ``will generally contain much less information than the raw training data'' \cite{pmlr-v54-mcmahan17a}.

\chapter{Introduction}\label{chap:intro}
\section{Context of the Thesis}\label{sec:context}

Machine Learning (ML) has become a central component of many digital services and supports decisions and predictions in many domains. At the core of these capabilities lies the ability of learning algorithms to extract statistical properties from data and use them to generalize beyond the observations available during training. In many learning settings, larger and more diverse datasets allow models to capture richer patterns and improve their predictive capabilities~\cite{banko-brill-2001-scaling,shalev2014understanding,hestness2017deeplearningscalingpredictable,kaplan2020scalinglawsneurallanguage}. The growing availability of digital data has therefore become one of the main resources enabling the development of increasingly capable machine learning systems.

Most of the information used to train these models originates from human activity and describes individuals, either directly or indirectly. A sequence of locations may reveal where someone lives or works; consumption and browsing histories may expose interests, habits, or preferences; medical and financial records describe particularly sensitive aspects of a person's life.
Even if harmless alone, when combinaed, these infotmation traces can become sensitive~\cite{netflix_attack}. From the perspective of a learning algorithm, individual records form a dataset from which population-level patterns can be inferred, but, from the perspective of the individual, each observation may correspond to a fragment of their private life.


For this reason, personal information is subject to legal requirements governing its collection and use, including principles such as purpose limitation and data minimization, as well as rights granted to the individuals concerned~\cite{eu_gdpr_2016,california_ccpa_cpra}. More recent regulatory frameworks further address the use of sensitive data and the deployment of machine learning systems in specific domains~\cite{eu_ai_act_2024,eu_ehds_2025,Buick_copyright}. Yet legal safeguards do not eliminate the underlying imbalance between the individuals who generate the data and the organizations that use them to provide digital services. Once data leave the direct control of their owners, those individuals lose control on how the information may later be exploited.

This imbalance is particularly significant in modern digital ecosystems, where a relatively small number of organizations operate many of the devices, applications, platforms, and cloud services through which everyday activities are mediated. Their position gives them access to detailed information about users' behavior, preferences, movements, purchases, communications, and social interactions. Such knowledge enables valuable and highly personalized services, but can also be used for profiling, targeting, ranking, persuasion, facial recognition, and other forms of automated decision-making that may remain opaque to the people affected~\cite{pasquale_black_box_society_2015,veliz_privacy_is_power_2021}. 


 Recent enforcement actions involving major technology companies demonstrates that privacy risks often arise not only from external attacks or accidental data breaches, but also from ordinary platform practices.
The Cambridge Analytica scandal provides a prominent example, where personal information collected through a digital platform was subsequently exploited for political profiling and targeted campaigning, raising concerns about consent, transparency, and secondary uses of personal data~\cite{ico_democracy_disrupted_2018,ico_data_analytics_political_campaigns_2018}. This episode is part of a broader pattern of documented privacy controversies involving major technology companie, including cases related to advertising practices, children's data, voice recordings, and the reuse of information collected for unauthorized purposes~\cite{cnil_google_gdpr_2019,ftc_youtube_coppa_2019,ftc_twitter_security_data_ads_2022,ftc_amazon_alexa_coppa_2023,texas_ag_google_privacy_2025}. 
Similar concerns arise in the data collection process for machine learning sysyems, where the demand for large and diverse training datasets creates strong incentives to collect and reuse data at scale. Controversies involving biometric systems trained on scraped or improperly obtained facial images \citeweb{noyb_clearview_criminal_complaint_2025}, copyrighted material used for model training \citeweb{reisner_ai_pirated_books_problem_2025}, and wearable devices capturing data from both users and bystanders \citeweb{svd_meta_ai_glasses_privacy_2026} demonstrate how the demand for data can blur the boundary between legitimate use and unauthorized collection or reuse.  More generally, conventional centralized machine learning requires data held by different individuals or organizations to be brought together before they can be jointly processed. Even when such collection is legitimate, concentrating raw information under the control of a single entity reinforces the dependence on that entity to handle it appropriately.

A natural way to reduce this dependence is to limit the amount of raw information that must be disclosed in the first place. Rather than moving data towards a central learning infrastructure, computation can be moved closer to where the data are generated or stored, allowing different parties to collaborate without pooling their records. The principle is particularly relevant when information is already distributed across personal devices, organizations, or institutions and cannot, or should not, be centrally collected. Federated Learning (FL) provides a framework for learning under these constraints~\cite{pmlr-v54-mcmahan17a}.


\section{Federated Learning} \label{sec:fed_learning}

Federated Learning enables multiple participants to train a shared machine learning model while keeping their datasets at their original locations. Each client trains a model copy locally on its own data and shares only the resulting model update, such as a gradient or updated model parameters~\cite{konecny2015federatedoptimizationdistributedoptimizationdatacenter, pmlr-v54-mcmahan17a}.
The main advantage and motivation behind  FL is its its natural adaptation to large-scale data collection settings. In applications involving mobile and edge devices, information is naturally produced across a large population of geographically distributed clients, while transferring and persistently storing all of it in a central infrastructure may be costly or impractical. 

 In the canonical FL formulation, a central server orchestrates the training process: it distributes the current model to selected clients, receives their locally computed updates, and aggregates them to obtain the next global model. A fully decentralized alternative removes the central coordinator and lets participants exchange information directly according to a communication graph, collectively updating their models through local interactions. Both settings avoid centralizing the underlying datasets, but differ in who observes the exchanged information and in the trust relationships created by the learning protocol. Throughout this thesis, we focus on the client-server formulation of FL~\cite{pmlr-v54-mcmahan17a}.

A second important aspect of FL is that it  limits the direct disclosure of personal information. Keeping records at their original location is naturally aligned with the principle of \emph{data minimization}, according to which personal information should be limited to what is necessary for a specified purpose~\cite{eu_gdpr_2016,eu_regulation_2018_1725,TRUONG2021102402}. This characteristic has made FL particularly attractive in privacy-sensitive domains such as healthcare and finance~\cite{Rieke2020,ZHANG2021106775,10.1145/3298981}. However, as we will later discuss more in details in Section~\ref{sec:privacy_in_fl}, ensuring data locality is not enough to guarantee privacy.
\subsection{Federated Averaging}

We now formalize the federated learning setting, following the original formulation of \citeA{pmlr-v54-mcmahan17a}. Let $\mathcal{U}=\{1,\dots,U\}$ denote a finite set of participants, commonly referred to as clients, users, or devices, coordinated by a central server. The objective is to train a single global machine learning model with parameters $\theta\in\mathbb{R}^d$. Each client $u\in\mathcal{U}$ holds a local dataset $\mathcal{D}_u=\{(\mathbf{x}_i^u, y_i^u)\}_{i=1}^{N_u}$, where $N_u = |\mathcal{D}_u|$ denotes the number of samples available at client $u$. Local datasets may differ both in size and in their underlying data distributions, reflecting the statistical heterogeneity that is typical of federated learning scenarios.

The goal of federated learning is to minimize a global empirical risk obtained by aggregating the local objectives of the clients:

\begin{equation}
\min_{\theta\in \mathbb{R}^d}
\mathcal{L}(\theta)
=
\sum_{u=1}^U p_u \mathcal{L}_u(\theta),
\end{equation}
where $\mathcal{L}_u(\theta)$ denotes the local objective of client $u$ evaluated on its local dataset $\mathcal{D}_u$, and $p_u$ is the aggregation weight associated with client $u$. A common choice is $p_u = 1/U$ for all $u$, which assigns the same weight to each client. Alternatively, one may set $p_u = \frac{|\mathcal{D}_u|}{|\mathcal{D}|}$, where
$\mathcal{D} = \bigcup_{u=1}^U \mathcal{D}_u$, so that clients are weighted proportionally to the size of their local datasets.

The local objective is typically written as an Empirical Risk Minimization (ERM) problem:
\begin{equation}
\mathcal{L}_u(\theta)
=
\frac{1}{N_u}
\sum_{i=1}^{N_u}
\ell(\theta;\mathbf{x}_u(i),y_u(i)),
\end{equation}
where $\ell(\theta;\mathbf{x}_u(i),y_u(i))$ is the loss incurred by the model with parameters $\theta$ on the $i$-th sample of client $u$.

The original algorithm proposed to solve this optimization problem is Federated Averaging (FedAvg) \cite[Algorithm~1]{pmlr-v54-mcmahan17a}, an iterative procedure performed over $T$ communication rounds. Its main steps are summarized in Algorithm~\ref{algo:fl-framework}. As discussed above, the training process is coordinated by a central server, which initializes the global model $\model$ (line 1). At each communication round $t$, the server selects a subset of participating clients, denoted by $\mathcal{U}^t \subseteq \mathcal{U}$ (line 3), and sends them the current global model (line 4). After receiving the model, each selected client performs local training on its private dataset (line 10), typically by running stochastic gradient descent for a fixed number of local epochs $E$ (Algorithm~\ref{algo:fedavg}). Once local training is completed, each client sends back an updated model $\model_u^{t+1}$ (line 5). The server then aggregates the received client models to obtain the next global model $\model^{t+1}$ (line 6).

FedAvg represents the canonical algorithmic formulation of federated learning. Since its introduction, many variants have been proposed, modifying different components of the protocol, including the client optimizer, the server aggregation rule, the client sampling strategy, and the communication pattern~\cite{reddi2021adaptivefederatedoptimization}. This line of work shows that federated learning should not be understood as a single fixed algorithm, but rather as a flexible framework for distributed optimization under data locality constraints. At the same time, this flexibility is needed because the federated setting departs from the assumptions of conventional centralized training. Clients may hold data drawn from different distributions~\cite{Li2020On, marfoq2021federated, wang_client, scaffold}, operate under different computational~\cite{li2019fedmdheterogenousfederatedlearning,FjORD} and network conditions \cite{caldas2019expandingreachfederatedlearning, chen_efficient, ALMANIFI2023100742}, participate intermittently \cite{ribeiro_client_corr, rodio_clients, sun_participation}, and have different privacy requirements \cite{brendan2018learning, Fu_dp_dl_review}. As a result, federated learning introduces a set of challenges that are either absent or less pronounced in centralized learning, including statistical and system heterogeneity \cite{ye_review_het_fl}, communication efficiency \cite{ALMANIFI2023100742}, fairness \cite{Vucinich_fairness, Shi_fairness}, personalization \cite{tan_personalization}, robustness \cite{Uddin_robustness}, and privacy \cite{fu_privacy}.

\begin{algorithm2e}[t]
\caption{FL Framework}
\label{algo:fl-framework}

\KwOutput{global model $\model^T$}

\textbf{Server $s$} initializes $\theta^1$\;
\For{$t \in \{0,\ldots,T-1\}$}{
    $s$ selects a subset of clients $\clients^t \subseteq \clients$\;
    \label{algoline:select_client}

    $s$ broadcasts the global model $\model^t$ to $\clients^t$\;
    \label{algoline:broadcast}

    $s$ waits for the updated models $\model_u^t$ from every client
    $u \in \clients^t$\;

    $s$ computes $\model^{t+1}$ by aggregating the received updated models\;
    \label{algoline:aggregation}
}

\BlankLine

\textbf{Client $u \in \clients$}\tcp*[f]{Input: $\model$; Output: $\model_u$}

\While{FL training is not completed}{
    $u$ listens for the arrival of a new global model $\model$\;

    $\model_u \leftarrow
    \localupdaterule^u(\model, \dataset_u)$\;
    \label{algoline:updaterule}

    $u$ sends $\model_u$ to the server\;
}
\end{algorithm2e}

\begin{algorithm2e}[ht]
  \caption{Client $u$'s local update rule in FedAvg}
  \label{algo:fedavg}

  \KwInput{server model $\model^t$; local dataset $\localdataset$;
           batch size $B$; local epochs $E$; learning rate $\eta$}
  \KwOutput{updated local model $\model_u^t(K)$}

  $\model_u^t(0) \leftarrow \model^t$\;
  $\mathcal{B} \leftarrow$ split $\localdataset$ into batches of size $B$\;
  $k \leftarrow 0$\;

  \For{each local epoch $e \in \{1,\ldots,E\}$}{
      \ForEach{batch $b \in \mathcal{B}$}{
          $\gradient(\model_u^t(k), b)
          \leftarrow
          \frac{1}{B}
          \sum_{x \in b} \nabla \ell(\model_u^t(k), x)$\;

          $\model_u^t(k+1)
          \leftarrow
          \model_u^t(k)
          - \eta \gradient(\model_u^t(k), b)$\;
          \label{algoline:gradient_update}

          $k \leftarrow k + 1$\;
      }
  }

  \Return{$\model_u^t(K)$, where $K = E\left\lceil \frac{S_u}{B} \right\rceil$}\;
  \end{algorithm2e}

\subsection{Challenges in Federated Learning}
\francesco{I am not sure if to keep this subsection or if to remove it, because it might sound out of scope}

\paragraph{Heterogeneity}
Unlike centralized training, where data are typically collected in a single infrastructure and processed under relatively uniform computational conditions, federated learning operates over a population of clients that contribute both data and computation to the training process. This creates two closely related forms of heterogeneity: \textit{statistical heterogeneity}, which concerns differences in the local data distributions, and \textit{system heterogeneity}, which concerns differences in the computational and communication capabilities of the clients~\cite{kairouz_open_problems}.


System heterogeneity arises from differences in clients' hardware, energy availability, memory, bandwidth, and availability over time. These factors determine whether a client can participate in a given communication round, how much local computation it can perform, and how quickly it can return its update to the server. Since such conditions are often outside the direct control of the server, client participation in federated learning is typically partial, intermittent, and uneven across the population.

This uneven participation can affect the training process in several ways. If some clients are more likely to participate than others, the aggregated updates may overrepresent specific regions of the client population, potentially introducing bias in the global model \cite{wang_client}. Moreover, temporal or spatial correlations in client availability can affect convergence and increase the variability of model performance across rounds. System heterogeneity also exacerbates the \textit{straggler effect}, where slower clients delay the completion of a training round or must be dropped from aggregation \cite{kairouz_open_problems}.


Statistical heterogeneity arises because each client generates data through its own local interactions, environment, behavior, and preferences. As a result, the local datasets may differ in size, feature distribution, label distribution, or more generally in the relationship between inputs and outputs \cite{kairouz_open_problems, ye_review_het_fl}. Consequently, the data available at each client may not be representative of the overall population, and the local objectives optimized by different clients may point toward different or even conflicting solutions~\cite{scaffold}. Training a single global model under these conditions can therefore lead to a model that performs unevenly across clients, or poorly represents clients whose data differ substantially from the majority~\cite{marfoq2021federated}. 

\paragraph{Personalization}
Closely related to statistical heterogeneity, personalization in federated learning addresses the problem that a single global model may not be equally suitable for all clients \cite{marfoq2021federated}. When clients have different data distributions, the model that performs best on average across the population may perform poorly for specific clients or subgroups. One possible solution would be for each client to train an entirely local model, without participating in the federated process. However, in many cases, individual clients do not have enough data to learn accurate and robust models in isolation. Conversely, participating in standard FL allows clients to benefit from the data indirectly contributed by others, but may produce a global model that fails to capture their local distribution.

Personalized federated learning aims to address this trade-off. The goal is to exploit shared information across clients while still adapting the learned model to each client's local data and task~\cite{kairouz_open_problems}. Several approaches have been proposed to achieve this balance. Some methods learn a common global representation and personalize only part of the model, such as the final layers or client-specific heads \cite{fedper, fedrep}. Others formulate personalization as a regularized optimization problem, where each client learns a local model constrained to remain close to a shared global model \cite{pfedme, ditto}. A further line of work uses clustering or mixture-based approaches, in which clients are grouped according to similarities in their data distributions or learned models, and each group is assigned a more specialized model \cite{marfoq2021federated, marfoq_memorization, fedGMM}
\paragraph{Fairness}
Personalization can improve local utility by adapting models to client-specific objectives and data distributions. Nevertheless, federated training should also ensure that this adaptation does not lead to systematically poor performance for under-represented clients, populations, or protected groups.
Fairness is therefore a central challenge in FL, and different deployments may require different fairness guarantees. Accordingly, several notions of fairness have been considered in the literature \cite{Vucinich_fairness,Shi_fairness, salazar_fairness}. For instance, one may require the final model to satisfy group fairness \cite{salazar_fairness}, similarly to centralized machine learning \cite{dwork_fairness, zemel2013learning_fairness}, by avoiding systematic disparities across protected groups defined by attributes such as race, gender, or age. Alternatively, FL introduces specific notions of fairness. A first one is client-level fairness \cite{li2020fairresourceallocationfederated, erraji2026lossgapparityfairness}, which requires the global model to provide comparable utility across clients, so that no client systematically benefits at the expense of others. A second one concerns participation in the training process. Since only a subset of clients is typically selected at each communication round, frequently selected clients may have a stronger influence on the learned model, whereas rarely selected clients may remain poorly represented \cite{huang_fairness, fairness_shi_client_selection}. Finally, fairness may also be understood in terms of contribution: clients that provide more useful updates, data, or computational resources may need to be compensated or incentivized proportionally to their role in the training process \cite{shi_fairness_incentives, Chen_fairness_incentives}. These different fairness notions raise distinct challenges and can interact with privacy in non-trivial ways \cite{chen_privacy_fairness, fairness_fl_rafi}. In particular, enforcing or auditing group fairness often requires access to sensitive attributes or group-level statistics, which may conflict with the privacy-preserving motivation of FL. As discussed in Chapter~\ref{chap:aia}, sensitive attributes used to define fairness groups may themselves be vulnerable to inference attacks during FL training.

\paragraph{Robustness} Machine learning models may be targeted by adversaries during both training and inference that aim to disrupt the normal model functioning. At training time, \textit{poisoning attacks} aim to corrupt the model by injecting malicious samples into the training set or perturbing existing ones \cite{biggio_poisoning, cina_poisoning, jagielski_poisoning}. Depending on the attacker’s objective, poisoning may broadly degrade model performance or introduce a \textit{backdoor}, causing an attacker-chosen prediction when a specific trigger is present while preserving normal behavior on clean inputs \cite{cina_poisoning}. At inference time, \textit{evasion attacks} instead manipulate test inputs to induce incorrect predictions, often through small but carefully crafted perturbations \cite{biggio_evasion, carlini_wagner_robustness}.

In federated learning, these threats are amplified by the distributed nature of the training process. Since the server aggregates updates produced by clients whose local data and computation are not directly observable, compromised participants can act as \textit{Byzantine} clients \cite{blanchard_byz}, i.e., clients that arbitrarily deviate from the protocol and submit malicious or corrupted updates. Such updates may aim to degrade the overall performance of the global model  \cite{pmlr-v97-bhagoji19a, alie_attack, foe_attack, fang_poisoning}, or implant a backdoor while preserving normal accuracy on clean validation data \cite{pmlr-v108-bagdasaryan20a, Xie2020DBA}. Compared to centralized poisoning, the attacker does not need to contaminate a centrally collected dataset; it can directly influence the optimization process through the updates sent to the server. This has motivated a large body of work on Byzantine-resilient aggregation \cite{blanchard_byz, mhamdi2018hiddenvulnerabilitydistributedlearning} and robust distributed optimization~\cite{yin2018byzantine, GuerraouiGP24, Allouah_bounds, Allouah_clipping} to protect and tolerate adversaries in the training process. 

\section{Privacy in Federated Learning}\label{sec:privacy_in_fl}
In Section~\ref{sec:fed_learning}, we introduced the main motivations behind federated learning, including the intuitive privacy benefit of keeping data local. Data locality, however, does not constitute a privacy guarantee. Although FL avoids transferring raw records to a central repository, the learning process still exposes information derived from local datasets through the trained model and the messages exchanged during training. Consequently, assessing privacy in FL requires understanding what an adversary can infer about clients or their data from the information available throughout the protocol.

In centralized machine learning, trained models may reveal information about the data used to train them~\cite{shokri_mia, aia_fredrikson, carlini_extract_data, feng2024privacy}. Federated learning introduces an additional attack surface because clients repeatedly communicate data-dependent model updates to the server. Such updates may reveal properties of individual records or local datasets even though the underlying data never leave the clients. Moreover, an adversary that actively participates in the protocol may manipulate the exchanged models or the training procedure to increase the amount of information revealed.

A large body of work has demonstrated that, without additional safeguards, such leakage can expose private or sensitive information about clients and their local datasets~\cite{lyu2020threatsfederatedlearningsurvey, zhao_fl_attacks_review}. The remainder of this section reviews the main classes of privacy attacks studied in federated learning, followed by the mechanisms proposed to mitigate them.

\subsection{Threat Models}\label{subsec:threatmodel}

Before reviewing the main classes of privacy attacks in federated learning, it is useful to clarify the assumptions made about the adversary. A threat model characterizes the adversary's position in the learning protocol, the information available to it, the actions it is allowed to perform, and the private information it aims to infer. A central distinction in FL is between \textit{honest-but-curious}, or passive, adversaries and \textit{malicious}, or active, adversaries~\cite{Nasr2018ComprehensivePA}.

In the honest-but-curious setting~\cite{honest_but_curious}, the adversary follows the prescribed training protocol while attempting to infer private information from the information exposed during its execution. A typical example is the FL server, which legitimately receives client updates and may analyze them to learn information about the underlying local datasets. Since the adversary does not interfere with the protocol, the attack does not introduce deviations in the training process that clients could directly detect.

A malicious adversary, in contrast, is allowed to deviate from the prescribed protocol to increase privacy leakage. Such adversaries may occupy different positions in an FL system, including compromised clients or a coalition of clients. In the privacy-attack literature, however, the central server is frequently considered as the adversary~\cite{Nasr2018ComprehensivePA, zhao_fl_attacks_review, carletti2025sok}. The server occupies a particularly powerful position because it orchestrates training, selects participating clients, and distributes the model on which local computation is performed. A malicious server can therefore modify the model parameters or other information sent to a targeted client so that the resulting update reveals more about its private data. Such attacks can achieve substantially stronger leakage than their passive counterparts, but rely on a stronger adversarial model and may become detectable when clients verify the received model or otherwise check that the server follows the expected protocol~\cite{mia_survey_bai, carletti2026detectability}.

The adversary's level of access provides a complementary characterization of the threat model. In a \textit{black-box} setting, the adversary interacts with a model only through its inputs and outputs, a setting commonly considered for attacks against already trained models. \textit{White-box} access gives the adversary knowledge of internal model information, such as its architecture and parameters, while \textit{grey-box} access describes intermediate forms of knowledge. The distinction is independent of whether the adversary is passive or active: an honest-but-curious FL server, for example, may have white-box knowledge of the model while still following the training protocol faithfully. Training-time privacy attacks in FL commonly assume extensive model knowledge because the server itself coordinates the learning process.

In this thesis, Chapter~\ref{chap:aia} considers both passive and active adversarial models, whereas Chapter~\ref{chap:ctp} and~\ref{chap:vgia} investigate active adversaries, with particular emphasis on malicious-server attacks.


\subsection{Membership Inference Attacks}

A core privacy question is whether an adversary can determine if a particular record was used to train a model. Even without recovering any of its attributes, revealing the presence of a record in the training set may disclose sensitive information about the corresponding individual, for example, in a medical domain, it can be used to infer if a patient was part of a control group or not~\cite{yeom2018privacyriskmachinelearning}. Membership inference attacks (MIAs) formalize this privacy threat by attempting to distinguish training records from records that were not used during learning \cite{shokri2017membershipinferenceattacksmachine,yeom2018privacyriskmachinelearning,mia_survey_bai}.

Formally, given a target training dataset $\dataset$, a candidate record $(\x,y)$, and an observation $\mathcal{O}$ available to the adversary, such as the final model, its outputs, gradients, model updates, or a sequence of intermediate models, a MIA attempts to estimate the record's membership status through a binary decision rule
\begin{equation}
\hat{m}
=
\mathcal{A}_{\mathrm{MIA}}\big((\x,y),\mathcal{O}\big)
\in \{0,1\},
\end{equation}
where $\hat{m}=1$ indicates that the adversary predicts that $(\x,y)$ belongs to $\dataset$, while $\hat{m}=0$ indicates predicted non-membership.
MIAs are widely used as an empirical tool for assessing privacy leakage, since successful inference indicates that the information available to the adversary allows training and non-training records to be distinguished. They have also been used to audit data deletion and machine-unlearning procedures and as a basis for attacks seeking more detailed information about the provenance of a record~\cite{xiao_mia,hu2021sourceinferenceattacksfederated}.

MIAs were first studied in centralized learning~\cite{shokri2017membershipinferenceattacksmachine,yeom2018privacyriskmachinelearning,carlini2022membershipinferenceattacksprinciples}, where the adversary typically interacts with a trained model and exploits differences in its behavior on training and non-training records, for example through prediction confidence, loss, or correctness. Federated learning exposes additional information during the training process: depending on its position in the protocol, an adversary may observe intermediate models or client updates in addition to the final model. Membership inference has consequently been studied under both passive and active FL threat models~\cite{mia_survey_bai}.

The first FL-specific MIA was introduced by \citeA{Nasr2018ComprehensivePA}, who proposed a white-box attack in both passive and active settings. Under threat model, the adversary is assumed to know the target model architecture, its parameters, and the hyperparameters needed to run inference. Therefore, for any candidate record $(\x,y)$, the adversary can compute not only the model prediction $f(\x;\theta)$, but also hidden-layer activations, the loss $\mathcal{L}(f(\x;\theta),y)$, and the gradients of this loss with respect to the parameters of each layer. In the supervised version of the attack, the adversary also has access to an auxiliary dataset $\dataset_{\mathrm{aux}}$, containing examples with known membership status. These examples provide the labels needed to train a discriminator model. The attack model then takes, for each candidate record, a collection of white-box features as input, including layer-wise gradients, hidden-layer outputs, the final prediction vector, the one-hot encoding of the true label, and the loss value, collected over multiple training rounds, and outputs the membership status. 


Related work further demonstrated that the information exchanged during collaborative training can reveal the presence of individual records and other properties of participants' local data~\cite{melis2018exploitingunintendedfeatureleakage}. Subsequent MIAs have sought to reduce the auxiliary information and computational cost required by early attacks while exploiting signals specific to the iterative nature of FL. Some approaches improve or replace the data used to train the attack model, for instance through generative adversarial networks~\cite{zhang_gan_mia}. Another line of work exploits how membership signals evolve across communication rounds, using the trajectories of model outputs~\cite{zari2021efficientpassivemembershipinference}, prediction confidence~\cite{gu_mia}, bias parameters~\cite{Zhang_mia}, or adversarial robustness~\cite{liu_mia}. Other approaches avoid shadow-model training entirely. \citeA{li2023effective} exploit the geometry of gradient updates across consecutive rounds, while \citeA{Zhu_mia} formulate membership inference as a hypothesis test based on observations collected from multiple clients and communication rounds.
% MIAs have also been adapted to application-specific FL settings; for instance, \citeA{yuan2023interaction} study interaction-level membership inference in federated recommender systems, where the goal is to determine whether a user interacted with a particular item.

A malicious adversary can additionally modify the FL protocol to deliberately amplify leakage signals. For instance, \citeA{Pichler_mia} modify the parameters of a ReLU-based fully connected layer so that the resulting gradients reveal whether a queried sample participated in training. \citeA{mia_nguyen} demonstrate that malicious model parameters can amplify membership leakage even when clients employ local differential privacy. A different strategy is considered by \citeA{He_mia}, who poison the training data to encourage the model to memorize information associated with specific target classes before performing membership inference.
% \paragraph{Source Inference Attacks}
% In federated learning, membership inference can be extended to identifying the client whose local dataset contains a target record. Assuming that each training record is associated with a unique client and that $\dataset=\bigcup_{u=1}^U \localdataset$, a source inference attack can be expressed as

% \[
%     \hat{m}=\mathcal{A}_{\mathrm{source}}((\x,y),\mathcal{O}) \in \{0,1,\ldots,U\},
% \]
% where $0$ denotes non-membership and $u \in \{1,\ldots,U\}$ indicates that $(\x,y)$ is inferred to belong to the local dataset $\dataset_u$ of client $u$. Source inference is particularly relevant in cross-silo FL: for example, when several hospitals collaboratively train a model, identifying the institution that contributed a record may itself reveal additional sensitive information about the corresponding individual~\cite{hu2021sourceinferenceattacksfederated}.
%  A related formulation is subject-level membership inference, where the adversary aims to determine whether the data of a given individual participated in the federation, possibly through multiple records distributed across several institutions \cite{suri2022subject}.
\subsection{Attribute Inference Attacks}\label{sec:aia}

While membership inference attacks reveal whether a record was used during training, attribute inference attacks (AIAs) address a complementary question: what sensitive information about a partially known record can be inferred from the model or from the training process? Let a data point be decomposed as $(\x^p,s,y)$, where $\x^p$ denotes the public attributes, $s\in\mathcal{S}$ is a sensitive attribute with a set $\mathcal{S}$ of possible values, and $y$ is the associated label or target value. Given the public attributes and target value, together with an observation $\mathcal{O}$ available to the adversary, an AIA estimates the sensitive attribute as
\begin{equation}
     \hat s = \mathcal{A}_{\mathrm{AIA}}(\mathcal{O}, \x^p, y) \in \mathcal{S} .
\end{equation}

The observation $\mathcal{O}$ depends on the threat model and may include, for example, the output of a trained model, its parameters, gradients, model updates, or intermediate models observed during training.

An early influential instance of attribute inference was introduced by~\citeA{aia_fredrikson} in the medical domain. The adversary combines access to a trained model with partial information about a target individual to infer a missing sensitive attribute. A common model-based formulation evaluates candidate values of the sensitive attribute and selects the one for which the completed record best agrees with the model. Given a trained model $\theta$, this can be expressed as
\begin{equation}\label{eq:aia_mb}
\hat{s}
\in
\argmin_{s\in\mathcal{S}}
\ell\bigl(\theta,(\x^p,s,y)\bigr).
\end{equation}

The intuition behind the attack is that the learned model may encode information about how the sensitive attribute affects the prediction on records similar to those seen during training. As a result, when the adversary completes a partially known record with the correct sensitive value, the resulting input may be more consistent with the model behavior than completions based on incorrect candidate values. Such leakage may be amplified by memorization, overfitting, or overlearning, as the model can retain information that goes beyond the general relationships needed for the prediction task~\cite{yeom2018privacyriskmachinelearning,Song2020Overlearning}. At the same time, sensitive attributes may already be predictable from correlations with the public attributes, even without access to the trained model. \citeA{Jayaraman_aia} therefore compare AIAs against imputation baselines, emphasizing that privacy leakage should be assessed through the additional information provided by model access, rather than through reconstruction accuracy alone.

Later work refined both the attack objectives and the adversarial assumptions. For classification models, \citeA{mehnaz2022sensitiveattributesprivatenovel} propose black-box model inversion AIAs that use confidence scores to select the most likely sensitive attribute value, and also introduce a label-only variant that requires only the predicted class. Their work further considers a relaxed setting in which some non-sensitive attributes of the target record are also unknown to the adversary.

In federated learning, attribute inference can additionally exploit the model updates exchanged during training, providing the adversary with information that is unavailable from the final model alone, as in the centralized setting. Existing FL-specific AIAs are predominantly gradient-based: given the public attributes and labels of the target client's records, the adversary treats the unknown sensitive attributes as optimization variables and searches for the values whose corresponding gradients best match the observed client update~\cite{lyu2021novelattributereconstructionattack,chen2024practical}. In this way, the sensitive attributes are inferred from their contribution to the gradients produced during local training.

AIAs have been studied under both honest-but-curious and active threat models. A passive adversary only exploits updates naturally exposed by the protocol, whereas an active server can manipulate the model sent to the targeted client to strengthen the victim's signal carried by subsequent updates. A different malicious-server strategy is proposed by~\citeA{Gomes_aia}, who modify selected neurons so that information about sensitive attributes is directly encoded in the resulting gradients.

Model-based AIAs, despite their earlier use in centralized learning, have remained largely unexplored in federated settings. Chapter~\ref{chap:aia} addresses this gaps by studying model-based attribute inference for federated regression tasks.
% In particular, it shows how an adversary can approximate the target client's local model and use it to infer sensitive attributes, overcoming limitations of gradient-based attacks in regression in both honest-but-curious and malicious settings. 



\subsection{Data Reconstruction Attacks}\label{sec:dra}

The attacks discussed so far assume that the adversary already has access either to the record whose membership is to be inferred (MIAs), or to a partially observed record (AIAs). Instead, Data Reconstruction Attacks (DRAs)---referred to also as Gradient Inversion Attacks (GIAs)--- consider a stronger form of privacy leakage. In this setting, the adversary aims to recover the training data themselves from information exposed by the learning procedure. In federated learning, this information typically consists of the gradients or model updates exchanged between clients and the server. Existing attacks can be categorized in three main families, namely \textit{optimization-based} , \textit{generative-based}, and \textit{analytical} attacks \cite{carletti2025sok, Du_gia_survey}. Optimization-based attacks can be further distinguished according to whether the reconstruction is performed directly in the input space or in a latent representation. The following sections review these families in more detail.

\subsubsection{Optimization-based attacks}\label{subsec:optim_gradient_inversion}
Optimization-based attacks reconstruct private data by solving an inverse problem. The adversary initializes a candidate reconstructed batch and optimizes its inputs, and possibly its labels, so that the gradients computed on this batch match the update observed during federated training. More formally, for a target client $u$ and a set of observed rounds $\mathcal{T}$, the reconstruction problem can be written as
    \begin{equation}\label{eq:gia_general}
        \{(\hat{\x}_i,\hat{y}_i)\}_{i=1}^{B}
        \in
        \argmin_{\{(\x_i,y_i)\}_{i=1}^{B}}
        \sum_{t\in\mathcal{T}}
        \Dist\left(
            \frac{\partial \ell\left(
                \model^t,
                \{(\x_i,y_i)\}_{i=1}^{B}
            \right)}{\partial \model^t},
            \model^t-\model_u^t
        \right),
    \end{equation}
    where $\Dist(\cdot,\cdot)$ denotes a distance or dissimilarity measure between the gradient computed on the reconstructed batch and the update observed by the adversary. 

The original attack formulation by~\citeA{dlg} instantiated this objective with a Euclidean matching loss. Subsequent work showed that the success of optimization-based GIAs depended not only on solving the inverse problem, but also on how the gradient-matching objective was defined. In particular, \citeA{geiping_ig} showed that image reconstruction could be substantially improved by matching the direction of the observed and reconstructed gradients through a cosine-based objective, and by adding total variation regularization~\cite{rudin_tv} to promote spatially coherent reconstructions and stabilize the optimization.

These early attacks, however, were mainly effective in relatively idealized settings. They often assumed that the target labels were known and achieved high-quality reconstructions only for very small batches consisting of few samples. In realistic settings, when a client computes gradients over a local batch, the update observed by the server contains information from all samples in that batch, rather than separate per-sample gradients. As the batch size grows, the signal associated with each individual record becomes increasingly difficult to disentangle. A first line of follow-up work therefore focused on making the attack setting more realistic by recovering unknown labels together with the inputs~\cite{zhaoIDLGImprovedDeep2020, maINSTANCEWISEBATCHLABEL2023, gengImprovedGradientInversion2024, yeHighFidelityGradientInversion2024, liTemporalGradientInversion2025a}. In parallel, other works improved reconstruction quality by enriching the optimization objective, either through additional regularization terms~\cite{yinSeeGradientsImage2021, hatamizadehGradViTGradientInversion2022,liE2EGIEndtoEndGradient2023, gengImprovedGradientInversion2024, Fan_fedleak} or through losses designed to better correlate with reconstruction quality~\cite{wangVarianceGradientAlso2022, hongForeseeingReconstructionQuality2024, yeHighFidelityGradientInversion2024, Fan_fedleak}.

Improving the objective alone is often insufficient in large-batch regimes, where the main difficulty is the separation of multiple sample contributions from a single observed update. Some works therefore modified the reconstruction algorithm itself to better handle this regime. For instance, \citeA{kariyappaCocktailPartyAttack2023} formulated gradient inversion as a blind source separation problem, aiming to separate the per-sample gradient contributions contained in the batch gradient. More recently, \citeA{liTemporalGradientInversion2025a} exploited information collected across different training steps and improved reconstruction by combining multiple candidate reconstructions of the same batch, drawing inspiration from robust aggregation techniques.

Another important limitation arises in realistic federated protocols such as FedAvg, where clients perform several local optimization steps before communicating with the server. In this setting, the exhanged update is no longer a single gradient, but the result of an entire local training trajectory. To address this mismatch, \citeA{dimitrovDataLeakageFederated2022} proposed to simulate the FedAvg update during reconstruction, so as to better reproduce the client-side computation and recover images from model updates. Later works approximated this local trajectory through surrogate models, using either linear~\cite{Zhu_sme_gia} or non-linear~\cite{xiaNonLinearTrajectoryModeling2025} trajectory models to improve reconstruction from multi-step updates.

Although most optimization-based GIAs were initially designed for images, the same general principle has also been adapted to other data modalities. In the text domain, \citeA{deng-etal-2021-tag-gradient} designed TAG, an attack specifically targeting Transformer-based language models. Later, \citeA{lamp_balunovic} proposed to guide sentence reconstruction with a language-model prior. These settings still benefit from the fact that reconstruction quality can often be assessed through visual or semantic inspection. In the tabular domain instead, this is not possible: records may contain continuous, categorical, or mixed features, and reconstruction quality may be difficult to assess without access to the ground truth. For this reason, \citeA{veroTabLeakTabularData2023} designed the only optimization-based GIA tailored to tabular data, together with an entropy-based heuristic for estimating reconstruction quality.

All the attacks reviewed so far assume an honest-but-curious adversary, who passively observes the updates exchanged during the normal execution of the FL protocol. A stronger threat model considers a malicious adversary, typically the server, who can actively modify the model parameters, the architecture, or parts of the training protocol to induce more informative client updates. In this setting, the attacker can modify the model parameters sent to the client so that the resulting update contains a stronger signal from selected training samples. For classification tasks, one strategy is to manipulate the classification layer to isolate or amplify the gradient contributions associated with a target class~\cite{GradFiltClasswiseTargeted, wenFishingUserData2022a}. The same principle can also be exploited across multiple communication rounds to progressively focus the leakage on specific samples~\cite{wenFishingUserData2022a}. Alternatively, \citeA{shan2025geminio} used a vision-language model to assist the adversary in crafting a malicious model that artificially increases the loss on target samples described through natural language, thereby making their contribution to the observed update more prominent.

\subsubsection{Generative-based attacks} 
\francesco{In this section I only wrote in prose how the attack works. I am not sure if to add the objective function that they minimize, to make more clear how generative models are employed in the optimization process}

The use of generative models in federated-learning privacy attacks evolved through several distinct paradigms. Early works employed GANs to approximate class-level distributions in victims’ private data. In the first generateive-based attack proposed by \citeA{hitajDeepModelsGAN2017a}, a malicious participant exploited and actively influenced the global model to train a GAN that generated prototypical samples of a target class available only to a victim participant. Instead, in the attack proposed by \citeA{wangInferringClassRepresentatives2018}, the server became the attacker, and used the client local updates as proxy for the discriminator. Subsequent studies shifted toward instance-level gradient inversion by constraining the reconstruction search with generative models. \citeA{jeon2021gradient} proposed GIAS, which replaced the direct optimization of the candidate input in Equation~\eqref{eq:gia_general} with an optimization over the output space of a pretrained generator. GIAS first optimized the latent code and then fine-tuned the generator parameters, thereby compensating for inaccuracies in the pretrained generative prior. For settings in which a suitable pretrained generator was unavailable, the authors further introduced GIML, a meta-learning approach that learned a generative prior across multiple gradient-inversion tasks constructed from updates observed during federated training. GGL~\cite{liAuditingPrivacyDefenses2022a} reconstructed private samples from defense-transformed gradients by searching the latent manifold of a pretrained GAN with adaptive gradient-free optimization. GIFD~\cite{fangGIFDGenerativeGradient2023a} extended this strategy beyond latent-space search by successively optimizing intermediate generator features. GI-SMN~\cite{qian_gismn} subsequently optimized the latent code of a pretrained StyleGAN-XL while regularizing the generated images through total-variation, $L_2$, and group regulariztions.

A complementary family of attacks replaced the fixed pretrained prior with a generator optimized during the inversion process~\cite{renGRNNGenerativeRegression2022a,xuCGIRConditionalGenerative2023,zhangGenerativeGradientInversion2023a}. Rather than optimizing the reconstructed images directly, these methods updated the generator parameters so that its outputs induced gradients matching the observed client update.

A more recent line of work shifted the reconstruction cost to an offline training phase by learning a reusable inversion model or generative decoder assuming access to auxiliary training data. Once trained, these models could be applied to newly intercepted updates without solving a separate gradient-matching problem for each reconstruction. LTI~\cite{wu2023learning} followed the most direct approach, training an inversion network on auxiliary sample–gradient pairs to map gradients to their corresponding inputs. Defensive transformations, such as noise addition, pruning, or compression, were incorporated when generating the training gradients, allowing the inversion model to circumvent defenses. With the same goal~\citeA{Yue_gia2023} designed an encoder-decoder inversion pipeline that solved the gradient optimization problem in a latent space, and then decoded the reconstruced image using a pretrained decoder.

FGLA~\cite{xueFastGenerationBasedGradient2023c} instead adopted a hybrid analytical–generative pipeline. It first separated the sample-level representations corresponding to the input of the classifier from the gradients of its fully connected layer, and then fed these representations to a trained generator that reconstructed the original inputs. Under a malicious-server threat model, Scale-MIA~\cite{shiScaleMIAScalableModel2025} trained a generative decoder on auxiliary data and crafted linear layers in the latent space of the distributed model. The attacker could then recover latent representations from aggregated updates through a closed-form procedure and map them back to the input space using the trained decoder. SEER~\cite{garovHidingPlainSight2024} instead jointly trained the shared model and a secret server-side decoder so that the contribution of individual samples could be separated from the aggregated client gradients in a hidden representation space and then decoded offline on the server.

More recent work explored diffusion models as an additional prior for gradient inversion. GGDM~\cite{guFederatedLearningVulnerabilities2024} and gradient-guided conditional diffusion methods~\cite{huangGradientGuidedConditionalDiffusion2024, mengEnhancedPrivacyLeakage2025} used leaked gradients to steer the reverse diffusion process toward reconstructions consistent with the observed update. GUIDE~\cite{carlettiGUIDEEnhancingGradient2025}, instead, employed diffusion after the inversion process to refine and denoise the reconstructed samples. Mjölnir~\cite{liu2025mjolnirbreakingshieldperturbationprotected} trained a diffusion model to approximate the original gradients from perturbed observations and subsequently applied a conventional optimization-based reconstruction attack to the recovered gradients.

Despite this progress, generative-based attacks remain limited in large-batch settings, where disentangling the contributions of multiple samples from an aggregated update remains challenging. High-quality reconstruction in such regimes often requires combining the generative prior with complementary mechanisms that separate individual gradient contributions or actively amplify the leakage of selected samples. Empirical evaluations have also remained predominantly focused on image data, leaving the effectiveness of generative-based attacks on other data modalities less established.

\subsubsection{Analytical attacks}
The analytical study of gradient inversion can be traced back
to~\citeA{phong_gia}, who first observed that the input to an affine
layer can be recovered directly from its weight and bias gradients when
these gradients are computed on a single training sample. This observation later became the basis of analytical leakage attacks against fully connected layers followed by ReLU activations. Indeed, ReLU naturally induces sparse gradients, since samples yielding
non-positive pre-activations contribute zero gradient to the corresponding neuron.

To illustrate this mechanism, consider a fully connected layer with
$N$ neurons, parameterized by a weight matrix
$\W^{a} \in \mathbb{R}^{N \times d}$ and a bias vector
$\b^{a} \in \mathbb{R}^{N}$. For an input
$\x \in \mathbb{R}^{d}$, the post-activation output of neuron $i$ is
\begin{equation}
z^{a}*{i}
=
\operatorname{ReLU}
\left(
\W^{a}*{i}\x+b^{a}*{i}
\right).
\end{equation}
Let $\loss$ denote the loss associated with sample
$(\x,y)$. By the chain rule, the corresponding weight and bias
gradients are
\begin{equation}
\frac{\partial \loss}{\partial \W^{a}*{i}}
=
\frac{\partial \loss}{\partial z^{a}*{i}}
\frac{\partial z^{a}*{i}}{\partial \W^{a}*{i}}
=
\frac{\partial \loss_j}{\partial z^{a}*{i}}
\x
\indicator_{z^{a}*{i}>0},
\end{equation}
and
\begin{equation}
\frac{\partial \loss_j}{\partial b^{a}*{i}}
=
\frac{\partial \loss_j}{\partial z^{a}*{i,j}}
\frac{\partial z^{a}*{i}}{\partial b^{a}*{i}}
=
\frac{\partial \loss_j}{\partial z^{a}*{i}}
\indicator_{z^{a}*{i}>0}.
\end{equation}
Thus, whenever neuron $i$ is active and its bias gradient is non-zero,
the input can be recovered exactly as
\begin{equation}
\x
=
\frac{\partial \loss}{\partial \W^{a}*{i}}
\left(
\frac{\partial \loss}{\partial b^{a}_{i}}
\right)^{-1}.
\end{equation}

If the attacked layer is the first layer of the network, this expression
reveals the original training sample, otherwise, it reveals the
intermediate representation supplied as input to that layer. The
reconstruction is obtained in closed form and does not require an
iterative optimization procedure. When gradients are aggregated over a
batch, however, the same ratio generally produces a weighted linear
combination of all samples activating the neuron. Disentangling these
contributions is therefore the central problem addressed by
analytical attacks.

In the honest-but-curious setting,
\citeA{zhuRGAPRecursiveGradient2021} extended closed-form inversion from
shallow affine models to deep fully connected and convolutional networks
through a recursive, layer-wise reconstruction procedure. They also
introduced a rank-based analysis to identify the layers and
architectures for which exact recovery is possible, although its
strongest results concern single-example gradients and networks
satisfying suitable rank and invertibility conditions.
\citeA{pan_gia} subsequently characterized batch reconstruction through
the notion of \emph{neuron exclusivity}. A ReLU neuron is exclusive to a
sample $\x$ when it is activated by $\x$ but remains inactive for
every other sample in the batch. In this case, its aggregated weight and
bias gradients contain only the contribution of $\x$, enabling its
closed-form recovery. \citeA{luAPRILFindingAchilles2022} showed that
analytical inversion can also be extended to self-attention layers and
that learnable positional embeddings may expose the ordering of the
reconstructed image patches. In a complementary theoretical setting,
\citeA{wang2023reconstructingtrainingdatamodel} established
identifiability guarantees for sufficiently wide shallow and deep
networks and proposed a tensor-decomposition method for reconstructing training samples from a gradient evaluated at suitably selected model parameters. These passive attacks are nevertheless limited by restrictive conditions, such as single-samples gradients, or natural occurrence of the exclusive activations. Then, SPEAR~\cite{dimitrovSPEARExactGradient2024} addresses the single-sample limitation by combining a low-rank factorization of fully connected-layer gradients with the sparsity induced by ReLU activations. Using dictionary learning techniques~\cite{pmlr-v23-spielman12}, it enables exact reconstruction of complete batches of up to approximately 25 high-dimensional images. However, its sampling procedure has exponential complexity in the batch size. \citeA{bakarsky2025spearscalinggradientinversion} later addressed this limitation by replacing the search procedure with more efficient dictionary-learning techniques. For language models, DAGER~\cite{petrovDAGERExactGradient2024} builds on SPEAR's low-rank intuition by exploiting the low-dimensional subspace spanned by batch gradients, which contains the private inputs.

Rather than relying on ReLU-induced sparsity to disaggregate continuous samples, DAGER uses the finite vocabulary and known embedding matrix to identify tokens contained in the gradient-induced subspace and reconstruct complete textual batches. \citeA{drencheva2025grain} later extended this theoretical framework to GNN layers, recovering both client input features and graph structure under the assumption of discrete input features.

A distinct line of work considers a malicious server that is allowed to manipulate the model distributed to clients. A seminal example is the Robbing the Fed attack~\cite{fowlRobbingFedDirectly2021}, which introduces an \emph{imprint module} consisting of a maliciously configured fully connected layer followed by a ReLU activation. The neurons share the same weight vector and therefore compute the same linear projection of the input, while their biases are set to different ordered thresholds derived from prior knowledge of the projection distribution. As a result, consecutive neurons activate on nested subsets of the batch, and subtracting their weight and bias gradients isolates the contribution of samples falling between consecutive thresholds. When a single sample falls in such an interval, its input can be reconstructed in closed form from the gradients of a single communication round.

The imprint module subsequently became a fundamental building block for several attacks. Decepticons~\cite{fowlDecepticonsCorruptedTransformers2022} combines an
imprint-style linear leakage module with malicious modifications of the
attention layers, allowing the attacker to associate recovered
tokens with their positions in a sentence and reconstruct the original sequences in
transformer-based language models.
Panning for Gold~\cite{chu2023panning} further adapts this construction
from indiscriminate reconstruction to the targeted extraction of
privacy-sensitive information. Modified attention heads tag sequences
containing an attacker-defined trigger, while the linear leakage module
filters out irrelevant tokens and exposes only the targeted
continuations. Scale-MIA~\cite{shiScaleMIAScalableModel2025} applies the same imprint
mechanism to an intermediate layer of the target model. Rather than
recovering the original inputs directly, the malicious linear layers
isolate and reconstruct their latent representations in closed form,
which are subsequently mapped back to the input space using a
generative decoder trained on auxiliary data.
\citeA{Zhao_cvf} examined the resource overhead of linear-layer leakage under Secure Aggregation~\cite{secure_aggregation_bonawitz}, where the server observes only the aggregate of the participating clients' updates. Prior attacks effectively treat this aggregate as the update of a single large batch, requiring the size of the malicious fully connected layer to scale with the total number of samples across all clients and leading to prohibitive model growth. They proposed a sparse construction that separates the scaling with the number of clients from that with the number of local samples. This line of work was subsequently extended by LOKI~\cite{zhaoLOKILargescaleData2023a}, which combines client-specific convolutional mappings with fully connected imprint layers. The convolutional mappings keep gradient contributions from different clients separated even after Secure Aggregation, while the imprint layers isolate individual samples within each client.

An alternative ReLU-leakage attack was proposed by~\citeA{boenischWhenCuriousAbandon2023} through the \emph{Trap Weights} attack. Standard weight initializations may naturally cause some neurons in a fully connected ReLU layer to activate for only one sample, thereby exposing
that sample through linear-layer leakage. Trap Weights amplifies this
phenomenon by inconspicuously rescaling the weights of the attack layer,
increasing the number of neurons that isolate individual samples without
modifying the model architecture. A theoretical bound on the
reconstruction capability of this attack is discussed in
Chapter~\ref{chap:ctp}. A related approach was proposed
by~\citeA{ZHANG2023119421}, who trained malicious parameters for an
existing fully connected layer using an auxiliary dataset. Their method
extends analytical extraction beyond ReLU networks to models employing
Tanh and Sigmoid activations. A different architecture-preserving strategy is adopted by
MKOR~\cite{wangMaximumKnowledgeOrthogonality2024}. The attack
maliciously modifies the parameters of existing convolutional and fully
connected layers to maximize the orthogonality between prior knowledge
about the input structure and the information revealed by the gradients.
For images, this prior is based on the assumption that neighbouring
pixels tend to have similar intensities. Under this assumption, MKOR
enables analytical reconstruction of multiple samples from large
batches without requiring auxiliary data or changes to the model
architecture.

Finally, \citeA{pasquiniEludingSecureAggregation2022a} showed that a malicious server can defeat the protection offered by Secure Aggregation through model inconsistency. In their gradient-suppression attack, the server sends the honest model to the target client and manipulated models to the remaining clients, causing the latter to return zero updates. The resulting aggregate therefore coincides with the target client's individual update, which can subsequently be exploited by reconstruction attacks.

Building on this line of work, Chapter~\ref{chap:ctp} introduces a
geometric perspective on analytical gradient leakage and proposes an
attack capable of recovering all samples in a batch. Chapter~\ref{chap:vgia} then
investigates the previously overlooked problem of verifiability and
introduces a method that certifies the exact isolation of a private
sample before reconstructing it analytically.




\subsection{Other Attacks}
Federated and decentralized learning are subject to several additional forms of privacy leakage related to the attacks discussed above. Source inference attacks extend the membership inference problem by seeking to identify which client contributed a known or suspected training record~\cite{hu2021sourceinferenceattacksfederated, Hu2024sia, ijcai2025p536}. Property inference attacks instead target characteristics of a client's local dataset rather than attributes of individual records, as in AIAs. They aim to infer statistical or semantic properties of the local data, including properties unrelated to the intended learning task. Examples include the presence or prevalence of particular demographic groups, visual concepts, or medical conditions~\cite{melis2018exploitingunintendedfeatureleakage, Zhang_pia, Wang_pia, Kerkouche_pia}. 
Privacy leakage can also arise in fully decentralized learning, where participants exchange information directly according to a communication graph. Adversaries may exploit their position in the network and the dynamics of gossip-based algorithms to reconstruct information held by neighbouring or even non-neighbouring participants~\cite{Pasquini_insecurity, ElMrini_graph}. A final, related line of work concerns empirical privacy auditing. Canary-based methods deliberately introduce crafted records, gradients, or client updates and subsequently test whether their presence can be detected from the resulting model or exchanged information. Since the injected signal is known to the auditor, such methods can be used to quantify membership leakage or derive empirical lower bounds on the privacy loss provided by differential privacy~\cite{maddock2023canife, Nasr_tight_auditing, andrew2024oneshot}. Collectively, these results illustrate the variety of information that may be exposed by distributed training. 

\subsection{Privacy Mechanisms in Federated Learning}\label{sec:privacy_mechanisms}
Several mechanisms have been proposed to mitigate privacy risks in federated learning. Privacy protection, however, cannot generally be reduced to the presence or absence of a single defense~\cite{kairouz_open_problems}. Different mechanisms address different sources of leakage and rely on different trust assumptions. Privacy-preserving FL systems therefore often rely on multiple complementary layers of protection, selected according to the information exposed during training and the capabilities assigned to the adversary.

This subsection provides an overview of the main privacy-preserving mechanisms used in federated learning. For more comprehensive discussions, the reader is referred to dedicated surveys on privacy and security in federated learning~\cite{Rao_privacy_survey, Yin_defense_taxonomy, sok_secure_aggregation}.

Differential privacy (DP)~\cite{dp_dwork,dwork_dp_book} is widely regarded as the de facto standard for providing formal and quantifiable privacy guarantees. DP limits the influence that an individual contribution to a dataset can have on the output of an algorithm, making outputs obtained from adjacent datasets statistically difficult to distinguish.

Differential privacy (DP)~\cite{dp_dwork,dwork_dp_book} is widely regarded as the de facto standard for providing formal and quantifiable privacy guarantees. DP limits the influence that an individual contribution to a dataset can have on the output of an algorithm, making outputs obtained from adjacent datasets statistically difficult to distinguish. More formally, a randomized mechanism $\mathcal{M}$ satisfies $(\varepsilon,\delta)$-differential privacy if, for every pair of adjacent datasets $\mathcal{D}$ and $\mathcal{D}'$, and every measurable set $S \subseteq \mathrm{Range}(\mathcal{M})$,

\begin{equation}\label{eq:dp}
    \mathrm{Pr}\left[ \mathcal{M}(\mathcal{D} \in S)\right] \leq e^{\varepsilon} \mathrm{Pr}\left[ \mathcal{M}(\mathcal{D}' \in S)\right] + \delta.
\end{equation}
The parameter $\varepsilon$ bounds the multiplicative change in the probability of any output between adjacent datasets, with smaller values corresponding to stronger privacy guarantees. The parameter $\delta$ allows a small probability that this bound does not hold and is therefore typically chosen to be negligible. A common mechanism for training neural networks with record-level DP is DP-SGD~\cite{Abadi_dp}. At each optimization step, DP-SGD clips per-example gradients to bound their sensitivity and adds calibrated Gaussian noise to their aggregate.

In federated learning, different DP deployment models exist, depending on where randomization is applied and which parties are trusted. Under Local Differential Privacy (LDP), each client randomizes its contribution before transmitting it to the server. The server therefore never observes the corresponding unperturbed contribution, and privacy does not rely on trusting it with individual updates. However, because privacy must be ensured before aggregation, LDP generally requires substantially more noise than mechanisms that protect the aggregate, potentially degrading convergence and model utility, particularly for high-dimensional updates or strong privacy requirements~\cite{geyer2018differentiallyprivatefederatedlearning,kairouz_open_problems}.

Under \emph{central differential privacy} (CDP), a trusted server or aggregator receives individual client updates, bounds their sensitivity through clipping, computes their aggregate, and adds calibrated noise before the aggregate is used or released. Because the noise is applied to the aggregate rather than independently to every contribution, CDP typically provides a more favorable privacy-utility trade-off than LDP. The guarantee, however, relies on trusting the aggregator, which observes the individual updates before the privacy mechanism is applied. It therefore does not protect clients against a malicious or compromised server~\cite{mcmahan_user_dp,kairouz_open_problems}.

To improve utility relative to LDP without exposing individual unperturbed updates, \emph{Distributed Differential Privacy} (DDP) combines client-side randomization with secure aggregation~\cite{kairouz_discrete_gaussian, agarwal2021the, chen2022poisson}. Secure aggregation is a cryptographic multiparty-computation primitive that allows the server to recover only an aggregate, typically the sum of the clients' updates, while hiding each individual contribution~\cite{secure_aggregation_bonawitz,sok_secure_aggregation}. In DDP, each client clips its update and contributes a share of the noise required to achieve a target DP guarantee. Although an individual contribution does not necessarily satisfy LDP, the securely aggregated sum can provide a central-DP-like guarantee because the server observes only the noisy aggregate and cannot separate the contributions of individual clients~\cite{kairouz_discrete_gaussian, agarwal2021the, chen2022poisson}.

Distributed DP introduces additional assumptions and system costs. Its guarantees depend on sufficient client participation and on bounds regarding dropout and collusion, while secure aggregation adds computational, communication, and coordination overheads~\cite{kairouz_discrete_gaussian, chen2022fundamentalpricesecureaggregation}. Moreover, secure aggregation hides individual updates but does not verify the model distributed by the server or the validity of client contributions. In some settings, a malicious server may therefore circumvent its protection through inconsistent or manipulated models, without breaking the underlying cryptographic protocol~\cite{pasquiniEludingSecureAggregation2022a,zhaoLOKILargescaleData2023a}.

Despite its formal guarantees, differential privacy is not a universal solution to privacy leakage. Strong privacy typically requires aggressive clipping and substantial noise, which can slow convergence and degrade model accuracy, particularly when high-dimensional updates are exchanged over many communication rounds. Conversely, limiting the perturbation to preserve utility results in weaker privacy guarantees and may leave significant empirical leakage~\cite{Jayaraman_privacy_evaluation,cummings2024advancingdifferentialprivacyfuture}. Chapter~\ref{chap:aia} further examines this trade-off in the context of attribute inference attacks.

Beyond mechanisms providing formal privacy guarantees, several empirical countermeasures have been proposed against specific privacy attacks~\cite{bai_survey_mai,carletti2025sok}. For membership and property inference, researchers have investigated techniques such as selective gradient sharing, model perturbation, and weight pruning~\cite{melis2018exploitingunintendedfeatureleakage,yang_perturbation,Stripelis_pruning}. The protection offered by such techniques, however, is generally tied to the attacks and threat models considered during their evaluation.

Gradient inversion attacks have motivated a broader range of defenses operating at different stages of the learning process. Data-level approaches transform or refine sensitive samples before the corresponding updates are computed~\cite{jeterOASISOffsettingActive2024,fanRefinerDataRefining}. Other defenses reduce the information exposed during training by perturbing or obfuscating intermediate representations and model updates~\cite{sunSoteriaProvableDefense2021,Xue_pruning,zhangCENSORDefenseGradient2025}. Complementary approaches modify the model architecture~\cite{scheligaPRECODEGenericModel2022} or the learning procedure itself~\cite{Wudefense2024,Ye_privacy2023} to make reconstruction more difficult. In the malicious-server setting, client-side detection mechanisms have also been proposed to detect suspicious model manipulations before local training begins~\cite{wangHearNoEvil2025,Carletti_2026}.

Unlike differential privacy, however, such defenses generally provide empirical and attack-dependent protection rather than universal privacy guarantees. Their effectiveness may depend on the adversary's knowledge and capabilities, as well as on the model architecture, training configuration, and evaluation setting. Adaptive attacks have, in particular, been shown to circumvent defenses based on gradient compression, pruning, stochasticity, and noise-based perturbation~\cite{Yue_gia2023,scheligaCombiningStochasticDefenses2025,liu2025mjolnirbreakingshieldperturbationprotected,mengEnhancedPrivacyLeakage2025,carletti2025sok}. These limitations motivate a deeper understanding of the information exposed by federated learning protocols and of the conditions under which such information can be exploited. 

\section{Summary of the Thesis Contribution}

The preceding sections have shown that federated learning changes the trust assumptions underlying machine learning without eliminating them. Although clients retain their data locally, the central server still occupies a privileged position in the training process, coordinating client participation, distributing model parameters, and receiving the resulting updates. Federated learning therefore shifts part of the trust placed in centralized data custody toward the behavior of the coordinating server. If this trust is misplaced, the server may exploit its visibility and control over the training process to infer information about clients' private data. 

Privacy attacks provide a concrete means of assessing the risks associated with trust. By demonstrating what an adversary can infer from the information exposed during training, they establish empirical lower bounds on achievable privacy leakage and provide demanding baselines for evaluating privacy-preserving mechanisms. Studying increasingly capable adversaries is therefore important for understanding which privacy guarantees arise from the learning protocol itself and which instead depend on assumptions about the behavior of the parties involved.


Motivated by these considerations, this thesis pursues two main objectives:
\begin{enumerate}
    \item To deepen the understanding of privacy leakage in federated learning
    from the perspective of the central server, characterizing what information
    can be inferred from clients' updates under both honest-but-curious and
    malicious-server threat models.

    \item To develop stronger, more scalable, and more reliable privacy attacks
    that improve the empirical auditing of federated learning systems and
    provide more informative baselines for the design and evaluation of
    privacy-preserving mechanisms.
\end{enumerate}

The contributions of this thesis are motivated by a constructive objective. Federated learning can meaningfully contribute to privacy protection by limiting the disclosure of raw data, but its guarantees should not rely solely on the server behaving as expected. Effective protection also depends on the design of the training protocol, the information exposed during its execution, and the safeguards used to constrain or verify the behavior of participating entities. By clarifying what a central server can infer under different threat models, this thesis aims to support more rigorous privacy assessments and better-informed design choices for federated learning systems. The remainder of this section summarizes the individual contributions of the thesis.


\subsection*{Contribution \#1: Attribute Inference Attacks for Federated Regression Tasks.}

\noindent\textbf{Problem:}
Attribute inference attacks; federated regression; empirical privacy auditing.

\medskip
\noindent\textbf{Research gap:}
Existing attribute inference attacks for federated learning have predominantly
focused on classification tasks and rely on matching the gradients generated by
candidate sensitive attributes with the updates observed during training.
Their effectiveness in regression settings had remained largely unexplored.
Moreover, model-based attribute inference attacks, originally developed for
centralized learning, had not been adapted to exploit the information exchanged
during federated training. As a result, the poor performance of existing
gradient-based attacks could incorrectly suggest that federated regression
provides stronger privacy protection.

\noindent\textbf{Main contributions:}
\begin{itemize}
    \item We establish a theoretical lower bound on the accuracy of model-based
    attribute inference attacks for least-squares regression, showing how the
    attack success depends on the local model's training error and on the
    influence of the sensitive attribute on the prediction.

    \item We introduce model-based attribute inference attacks for federated
    regression under both honest-but-curious and malicious-server threat
    models. The proposed methods first approximate the optimal local model of
    the targeted client and then use it to infer the sensitive attributes of
    individual records.

    \item Through experiments on real-world datasets, we show that the proposed
    attacks substantially outperform existing gradient-based approaches,
    particularly under heterogeneous client distributions, and provide stronger
    empirical baselines for evaluating privacy leakage in federated regression.
\end{itemize}

\medskip
\noindent\textbf{Publication:}
 \textbf{Francesco Diana}, Othmane Marfoq, Chuan Xu, Giovanni Neglia,
    Frédéric Giroire, and Eoin Thomas.
    ``Attribute Inference Attacks for Federated Regression Tasks.''
    In \emph{Proceedings of the AAAI Conference on Artificial Intelligence},
    vol.~39, no.~15, pp.~16271--16279, 2025.



\subsection*{Contribution \#2: A Hyperplane-Based
Data Reconstruction Attack in Federated Learning.}

\noindent\textbf{Problem:} Gradient inversion attacks; malicious federated-learning servers; reconstruction
of large and low-dimensional data batches.

\medskip
\noindent\textbf{Research gap:}
Existing gradient inversion attacks either rely on auxiliary knowledge about the
clients' data distribution or lose effectiveness as the number of records
contributing to an update increases. In particular, sparsity-based analytical
attacks attempt to configure neurons so that individual records can be separated
from the remaining batch. Although effective for some high-dimensional image
data, their performance degrades considerably for large batches and
low-dimensional tabular records. Consequently, prior attacks could not reliably
recover all the records contained in large aggregated updates.

\noindent\textbf{Main contributions:}
\begin{itemize}
    \item We provide a geometric interpretation of existing sparsity-based
    analytical attacks and derive theoretical upper bounds on their
    reconstruction capabilities, explaining their limited effectiveness for
    large batches and low-dimensional data.

    \item We introduce a hyperplane-based reconstruction attack in which a
    malicious server controls the contribution of client records to the
    observed gradients and progressively isolates them across communication
    rounds, without requiring auxiliary data or prior knowledge of the client
    data distribution.

    \item We show theoretically and empirically that the proposed attack can
    recover complete batches independently of their size, given sufficiently
    many communication rounds, and demonstrate perfect reconstruction of image
    and tabular batches containing up to \(4{,}096\) records.
\end{itemize}

\medskip
\noindent\textbf{Publication:}
 \textbf{Francesco Diana}, André Nusser, Chuan Xu, and Giovanni Neglia.
    ``Cutting Through Privacy: A Hyperplane-Based Data Reconstruction Attack
    in Federated Learning.''
    In \emph{The 41st Conference on Uncertainty in Artificial Intelligence},
    2025.


\subsection*{Contribution \#3: No More Guessing: A Verifiable Gradient
Inversion Attack in Federated Learning}

\noindent\textbf{Problem:}
Verifiability of gradient inversion attacks; certified record reconstruction;
privacy auditing without access to ground-truth data.

\medskip
\noindent\textbf{Research gap:}
Gradient inversion attacks are generally evaluated by comparing reconstructed
records with the original client data. Such an evaluation is available to the
researcher conducting an experiment, but not to an actual adversary, who does
not know the private batch. Existing attacks may therefore not distinguish
correct records from incorrect ones. Visual inspection can sometimes provide a
heuristic assessment for images or text, but it is costly for large batches and
is largely ineffective for numerical tabular data. Moreover, the geometric
search introduced in the previous contribution relies on a manually selected
stopping threshold and cannot determine whether a non-empty region contains one
record or several records.

\medskip
\noindent\textbf{Main contributions:}
\begin{itemize}
    \item We introduce a verifiable gradient inversion attack that uses an
    algebraic subspace test to certify when records have been individually
    isolated, providing an intrinsic guarantee of reconstruction correctness
    without requiring access to the original client data.

    \item We develop an adaptive hyperplane-search procedure that analytically
    reconstructs certified inputs, recovers their corresponding targets, and
    reallocates subsequent queries to unresolved regions, thereby avoiding
    heuristic stopping thresholds and unnecessary communication rounds.

    \item We validate the attack on tabular and image datasets for both
    regression and classification tasks, demonstrating exact and automatically
    certified recovery under large-batch aggregation in settings where previous
    attacks either fail or cannot assess reconstruction correctness.
\end{itemize}

\medskip
\noindent\textbf{Publication:}
    \textbf{Francesco Diana}, Chuan Xu, André Nusser, and Giovanni Neglia.
    ``No More Guessing: A Verifiable Gradient Inversion Attack in Federated
    Learning.''
    \emph{arXiv preprint arXiv:2604.15063}, 2026.

\section{Publications}
\subsection*{Published}
\noindent\begin{itemize}
    \item  \textbf{Francesco Diana}, Othmane Marfoq, Chuan Xu, Giovanni Neglia,
Frédéric Giroire, and Eoin Thomas.
``Attribute Inference Attacks for Federated Regression Tasks.''
In \emph{Proceedings of the AAAI Conference on Artificial Intelligence},
vol.~39, no.~15, pp.~16271--16279, 2025.
    \item \textbf{Francesco Diana}, André Nusser, Chuan Xu, and Giovanni Neglia.
``Cutting Through Privacy: A Hyperplane-Based Data Reconstruction Attack
in Federated Learning.''
In \emph{The 41st Conference on Uncertainty in Artificial Intelligence},
2025.
\end{itemize}

\subsection*{Under Review}
\noindent\begin{itemize}
    \item \textbf{Francesco Diana}, Chuan Xu, André Nusser, and Giovanni Neglia.
    ``No More Guessing: A Verifiable Gradient Inversion Attack in Federated
    Learning.''
    \emph{arXiv preprint arXiv:2604.15063}, 2026.
\end{itemize}

\section{Thesis Outline}
The remainder of this thesis is organized as follows. Chapter~\ref{chap:aia}
studies attribute inference attacks for federated regression tasks. It introduces
model-based attacks for both honest-but-curious and malicious adversaries,
provides theoretical guarantees for least-squares regression, and evaluates
their effectiveness against existing gradient-based approaches.
Chapter~\ref{chap:ctp} turns to the complete reconstruction of clients' private
records. It provides a geometric analysis of the limitations of existing
analytical attacks and introduces a hyperplane-based method capable of
reconstructing large batches of image and tabular data.
Chapter~\ref{chap:vgia} addresses the problem of reconstruction verifiability
and proposes an attack that certifies when a private record has been correctly
isolated and recovered, while adaptively reducing the number of communication
rounds required by the reconstruction process. Finally, Chapter~\ref{chap:conclusion}
summarizes the main findings of the thesis and discusses possible directions and challenges
for future research on privacy attacks in federated
learning.

  % % !TEX root = ../sommaire.tex

\chapter{Introduction}

Cette classe \LaTeX\ est basée sur la classe \mytt{these-LINA} écrite par Frédéric Goualard. La classe a été modifiée pour correspondre au modèle de thèse demandé par l'ED STIC de Université Côte d'Azur.

Pour les autres ED, il suffit de supprimer le fichier \mytt{visuel.jpg} dans le dossier \mytt{img/} et de renommer le visuel correspondant à l'ED en \mytt{visuel.jpg}.

\section{Organisation}

  Pour compiler ce document vous pouvez utiliser le \mytt{Makefile}. Sinon, on m'a dit dans l'oreillette qu'il existait des éditeurs tel que Texmaker qui gèrent eux même la compilation.

  \subsection{Fichier principal et informations}
  Le fichier principal est le fichier \mytt{sommaire.tex}, la classe du document est \mytt{these-ISSS}.
  Différentes options sont possibles :
  \begin{itemize}
    \item \mytt{draft} pour une version brouillon du manuscrit, les liens sont désactivés et en noir, et les étiquettes des références sont affichées ;
    \item \mytt{revision} qui ajoute à chaque première page de chapitre la date et l'heure de la dernière compilation en pied de page ;
    \item \mytt{gray} pour une version en grayscale du manuscrit ;
    \item \mytt{final} (par défaut) pour la version finale du manuscrit ;
    \item \mytt{oneside} pour une impression recto et \mytt{twoside} (par défaut) pour une impression recto-verso ;
    \item \mytt{notitlepage} pour un manuscrit sans page de titre et \mytt{titlepage} (par défaut) avec page de titre.
  \end{itemize}

  Le titre du document ainsi que les informations correspondant au laboratoire, une co-tutelle ou les co-financeurs se trouvent dans le fichier \texttt{titreEtInfos.tex}.

  Les informations relatives au jury, au directeur, co-directeur et co-encadrant se trouvent dans le fichier \mytt{jury.tex}.

  Le résumée et les mots-clés en français et en anglais se trouvent dans le fichier \mytt{resume.tex}.

  \subsection{Les chapitres}
  À l'exception des remerciements et des notations, chaque chapitre a son dossier.
  Ce n'est pas nécessaire mais c'est plus facile pour s'y retrouver après, et dans chaque dossier il y a le fichier principal, et il peut y avoir d'autres fichiers ou dossiers (par exemple un dossier \mytt{img/}).

  Les chapitres sont ajoutés au document principal à l'aide des macros \verb|\input| ou \verb|\import|. À noter que la macro \verb|\include| ajoute forcément une nouvelle page et ne permet pas des inclusions en cascade, son utilisation est à éviter.

  \subsection{Les paquets}
  Les ajouts de paquets sont dans le fichier \mytt{mesmacros.sty}. Vous pouvez supprimer ou ajouter dans ce fichier d'autres paquets.

  Ce fichier utilise les fichiers suivant :
  \begin{itemize}
    \item \mytt{francisation.sty} dans lequel se trouvent les traductions des différents noms de section ou listes ;
    \item \mytt{couleurs.sty} dans lequel se trouvent des couleurs, et dans lequel vous pouvez ajouter des couleurs ;
    \item \mytt{theoremNames.sty} dans lequel se trouvent les définitions de théorèmes ;
    \item \mytt{theoremList.sty} dans lequel se trouve le code pour générer des listes de théorèmes, utilisé pour générer la liste des définitions et la liste des exemples à la fin de ce document ;
    \item \mytt{algo.sty} pour les algorithmes en français avec le package \mytt{algorithmic} et \mytt{colorationSyntaxique.sty} pour la coloration syntaxique des algorithmes avec le package \mytt{lstlistings}.
  \end{itemize}

  \subsection{PDF/A}

  Le document est désormais compilé pour êre compatible avec la norme PDF/A-3B. 
  Les modifications effectuées ont été faites à partir des instructions trouvées ici : \url{https://www.mathstat.dal.ca/~selinger/pdfa/}.

  Si cela ne vous convient pas vous pouvez commenter ou modifier les lignes 82 du fichier \mytt{mesmacros.sty} :

  \begin{framed}\vspace{-0.75em}
    \begin{verbatim}
\usepackage[a-3b]{pdfx}\end{verbatim}\vspace{-0.75em}
  \end{framed}

  Si vous voulez conserver cette norme veillez à modifier les informations contenues dans le fichier \mytt{sommaire.xmpdata}.

\section{Citations}

Dans ce document il y a 3 bibliographies :
\begin{itemize}
  \item la bibliographie pour les publications personnelles,
  \item la bibliographie pour les pages web,
  \item et la bibliographie générale.
\end{itemize}

Elles sont définies au début du fichier \mytt{sommaire.tex}.

\begin{framed}\vspace{-0.75em}
\begin{verbatim}
% Biblio pour les pages webs
\newcites{web}{Pages web}
% Biblio pour mes publications
\newcites{mine}{Mes publications}
\end{verbatim}\vspace{-0.75em}
\end{framed}

On utilise pour cela le package \mytt{multibib}, vous pouvez en ajouter ou renommer celles existantes, il ne faut pas oublier de modifier le fichier \mytt{Makefile} pour compiler les nouvelles bibliographies ou supprimer la compilation de celles que vous n'utilisez pas.

Par défaut, les bibliographies utilisent le package qui cite correctement des pages web.
Le package impose le style \mytt{apacite} qui utilise un système auteur-date.
Pour changer de style de la bibliographie, vous devrez probablement supprimer le package.

\paragraph{Les publications dans la bibliographie générale} doivent être citées avec la macro \mytt{\textbackslash cite\{key\}}.

\begin{framed}
\noindent\cite{T2012}\vspace{-0.75em}
\begin{verbatim}\cite{T2012}\end{verbatim}\vspace{-0.75em}
\end{framed}

\paragraph{Les publications pour la section pages web} doivent être citées avec la macro \mytt{\textbackslash citeweb\{key\}}.

\begin{framed}
\noindent\citeweb{AC9} \vspace{-0.75em}
\begin{verbatim}\citeweb{AC9}\end{verbatim}\vspace{-0.75em}
\end{framed}

\paragraph{Les publications pour la section correspondant à vos publications} doivent être citées avec la macro \mytt{\textbackslash citemine\{key\}}.

\begin{framed}
\noindent\citemine{DHHS2011} \vspace{-0.75em}
\begin{verbatim}\citemine{DHHS2011}\end{verbatim}\vspace{-0.75em}
\end{framed}

Pour référencer vos publications par la suite, vous pouvez utiliser la macro \mytt{\textbackslash cite} ou la macro \mytt{\textbackslash citemine}.
Si vous souhaitez que votre publication apparaisse dans la bibliographie générale à la fin du manuscript vous pouvez soit la citer avec \mytt{\textbackslash cite} ou l'insérer avec \mytt{\textbackslash nocite}. \nocite{DHHS2011}

Attention si vous utilisez un style de bibliographie numérotant les références, vous aurez des problèmes de numérotation si une publication est à la fois citée en utilisant la macro \mytt{\textbackslash cite} et la macro \mytt{\textbackslash citemine}.



Chaque bibliographie peut ensuite être affichée en utilisant les lignes suivantes :

\begin{framed}\vspace{-0.75em}
\begin{verbatim}\bibliographystyle<s>{apacite}
\bibliography<s>{biblio}
\end{verbatim}\vspace{-0.75em}
\end{framed}

où \verb|<s>| correspond à la bibliographie que l'on veut afficher.

Pour redéfinir l'affichage, il suffit de redéfinir le \verb|\bibtitle|, pour que ce soit un chapitre utilisez la commande suivante :

\begin{framed}\vspace{-0.75em}
\begin{verbatim}\renewcommand{\bibtitle}{\bibchapter}
\end{verbatim}\vspace{-0.75em}
\end{framed}

Pour que ça apparaisse comme une section, utilisez la commande suivante :

\begin{framed}\vspace{-0.75em}
\begin{verbatim}\renewcommand{\bibtitle}{\section*{\refname}}
\end{verbatim}\vspace{-0.75em}
\end{framed}


\section{Exemples d'utilisation}

Des exemples d'utilisation des différents paquets sont présents dans ce document.


Des citations pour compléter les bibliographies \citemine{HHSD2011,HHSD2012}, \citeweb{AC9}.


Des citations pour compléter les bibliographies \citeA{HHSD2011}.

\citeAweb{AC9}.

Une référence citée avec \mytt{\textbackslash citemine} \citemine{DHHS2011} et \mytt{\textbackslash cite} \cite{DHHS2011} pour voir.


\renewcommand{\bibtitle}{\section*{\refname}}
\bibliographystylemine{apacite}
\bibliographymine{biblio}

  % \input{notations}

  % \Part{État de l'art}
  \import{chapitre2/}{chapter2}
  \import{chapitre3/}{chapter3}

  % \Part{Contributions}
  \import{chapitre4/}{chapitre4}
  \import{chapitre5/}{chapitre5}

  \input{conclusion/conclusion}

  \backmatter

  \renewcommand{\bibtitle}{\bibchapter}
  \bibliographystyle{apacite}
  \bibliography{biblio}

  \renewcommand{\bibtitle}{\section*{\refname}}
  \bibliographystyleweb{apacite}
  \bibliographyweb{biblio}

  \listoffigures
%. \listoftables

  \renewcommand{\listtheoremname}{Liste des définitions}
  \listoftheorems[ignoreall,show={definition}]

  \renewcommand{\listtheoremname}{Liste des exemples}
  \listoftheorems[ignoreall,show={example}]

  \listofalgorithms


  \appendix
  \import{appendix/}{appendix_aia}
  \import{appendix/}{appendix_ctp}

\end{document}
